a) Bentuk diferensial tarik baliknya adalah

\mavergleichskettealign
{\vergleichskettealign
{\varphi^* ( xdx -zdy +dz )
}
{ =} { u^2 d(u^2)    -(-u+v^2) d(uv) +d (-u+v^2)
}
{ =} { 2u^3 du + (u-v^2) (vdu +udv) -du +2vdv
}
{ =} { (2u^3 +uv -v^3 -1) du    + (u^2 -uv^2 +2v )dv
}
{ } {
}
}
{}
{}{.}

b) Integral lintasannya adalah

\mavergleichskettealign
{\vergleichskettealign
{ \int_1^2 ( 2t^3 +1 -t^{-3} -1 )dt  +(t^2 -t^{-1} +2t^{-1} ) d t^{-1}
}
{ =} { \int_1^2 ( 2t^3 -t^{-3} ) dt  +(t^2 + t^{-1} ) d t^{-1}
}
{ =} { \int_1^2 ( 2t^3 -t^{-3} ) dt  -(1 + t^{-3} ) dt
}
{ =} { \int_1^2 ( 2t^3 - 2t^{-3} -1 ) dt
}
{ =} { \left(  { \frac{   1 }{  2 } }  t^4  + t^{-2}  -t  \right) \vert_{  1 }^{  2 }
}
}
{
\vergleichskettefortsetzungalign
{ =} { 8 + { \frac{   1 }{  4 } }  -2 - { \frac{   1 }{  2 } } -1 +1
}
{ =} { 5 + { \frac{   3 }{  4 } }
}
{ } {}
{ } {}
}
{}{.}

c) Lintasan komposisinya adalah

\mavergleichskettedisp
{\vergleichskette
{ \psi (t)
}
{ =} { (t^2,1,-t +t^{-2} )
}
{ } {
}
{ } {
}
{ } {
}
}
{}{}{.}
Dengan demikian,

\mavergleichskettealign
{\vergleichskettealign
{ \int_1^2 \psi^* (xdx -zdy +dz )
}
{ =} { \int_1^2 t^2d(t^2) + d (-t + t^{-2} )
}
{ =} {  \int_1^2 (2 t^3-1 -2t^{-3} ) dt
}
{ =} { 5 + { \frac{   3 }{  4 } }
}
{ } {
}
}
{}
{}{.}

 [[Kategorie:Latexseite]]
